\chapter{Présentation de l'entreprise AGT}
\label{chap:entreprise}

\agtlettrine{C}e chapitre présente la structure de AGT du stage à
l'origine de ce rapport. Il décrit d'abord l'activité et l'organisation
d'AG Technologies, puis l'environnement technique dont disposait
l'entreprise avant le développement d'AGT Infra, environnement dont les
limites ont directement motivé le projet détaillé dans les chapitres
suivants.
\minisommaire

\section{AG Technologies : activité et organisation}

Basée à Yaoundé, au Cameroun, AG Technologies est une entreprise
spécialisée dans la conception de solutions digitales, avec un
positionnement marqué sur les plateformes SaaS (Software as a Service),
l'intelligence artificielle et l'accompagnement à la transformation
numérique des entreprises africaines. Son offre se structure autour de
trois axes complémentaires : la conception de solutions digitales sur
mesure adaptées aux besoins spécifiques de chaque client, la prestation de
services techniques incluant le consulting stratégique et l'intégration de
systèmes complexes, et le développement de plateformes SaaS regroupant
applications métiers, outils cloud et solutions d'automatisation.

Sur le plan technique, l'entreprise couvre un périmètre large : développement
web et mobile (React, Next.js, Flutter et développement natif), solutions
cloud et SaaS (AWS, Azure, architectures scalables, pratiques DevOps),
intelligence artificielle (modèles de langage, vision par ordinateur,
traitement du langage naturel, apprentissage automatique), automatisation
de processus (agents conversationnels, agents vocaux, orchestration de
flux de travail) ainsi que conception d'API et de systèmes distribués
(microservices, REST, GraphQL, gRPC). Selon les éléments communiqués par
l'entreprise elle-même, cette activité s'est traduite par plus de 150
projets livrés et plus de 80 clients accompagnés, avec une présence
touchant plusieurs pays du continent.

C'est au sein de cette structure, et plus précisément de son équipe
technique, que s'inscrit le stage à l'origine de ce rapport.

\section{Environnement technique avant AGT Infra}

Avant le développement d'AGT Infra, AG Technologies exploitait un parc de
plusieurs serveurs sans outil centralisé de supervision ou de gestion. Une
partie de ces ressources reposait notamment sur des VPS (Virtual Private
Server) mutualisés, hébergeant plusieurs applications de l'entreprise tels que : AGT-BOT,
team-tool et AGT TaskFlow. La cohabitation de plusieurs applications imposait une vigilance particulière lors de toute
intervention, un incident sur l'une des applications pouvant affecter les
autres.

La gestion de ce parc reposait sur des connexions SSH manuelles, effectuées
au cas par cas selon les besoins : déploiement d'une mise à jour,
vérification de l'état d'un service, consultation de journaux applicatifs.
Cette approche ne permettait ni vue d'ensemble sur l'état du parc, ni
contrôle d'accès différencié selon les personnes intervenant sur les
serveurs, ni traçabilité centralisée des actions effectuées. C'est ce
constat, développé au chapitre suivant, qui a motivé la définition du
périmètre fonctionnel d'AGT Infra.